%--------------------------------------------------------------------------------
%
%--------------------------------------------------------------------------------
\unnumberedChapter{Screens}

The UI Elements library introduced in the previous chapter provides the
individual building blocks required for implementing a user interface. Buttons,
LEDs, rotary encoders and displays each encapsulate a particular aspect of the
interaction between the user and the hardware. While these elements are useful
individually, larger applications require a mechanism to organize them into a
coherent user interface.

Most embedded devices present their functionality through a sequence of screens.
Only a limited amount of information can be displayed at any one time, while the
available input devices are usually restricted to only a few buttons and perhaps
a rotary encoder. Consequently, user interfaces are commonly organized as a
hierarchy of screens between which the user navigates.

The purpose of the \textbf{UI Screen} library is to provide a lightweight
framework for implementing such user interfaces. Instead of managing button
events, display updates and menu navigation throughout the application, these
tasks are centralized in a simple screen hierarchy. Applications merely define
their individual screens and implement the functionality associated with them.

The library deliberately avoids imposing a particular visual appearance. It
provides the navigation framework and event dispatching while leaving the
presentation of each screen entirely to the application.

%--------------------------------------------------------------------------------
\unnumberedSection{Concept}

A user interface consists of a number of individual screens. At any given time,
exactly one screen is active and receives all user interaction events. Buttons,
rotary encoders and other UI elements therefore do not communicate directly with
the application. Instead, their events are forwarded to the currently active
screen.

Each screen is implemented as an object derived from the \texttt{UIScreen} base
class. The base class defines a set of virtual callback methods that are invoked
whenever user interaction occurs. Applications implement the behaviour of a
screen by overriding the callbacks relevant for the particular application.

A screen represents a display context. It controls the information presented on
the display and reacts to user input while it is active. Screens can be arranged
in a hierarchical structure. With the exception of the root screen, every screen
has a parent screen and may contain a list of child screens. This allows complex
user interfaces to be constructed from simple individual screens.

The navigation concept is intentionally simple. Two buttons are reserved as the
default navigation controls: \textbf{MENU} and \textbf{SELECT}. The MENU button
is normally used to move through the list of child screens. The list is
organized as a circular list, meaning that after the last child screen the next
MENU action returns to the first child screen.

The SELECT button activates the currently selected screen and therefore enters
its child hierarchy. When a new screen becomes active, the leaving screen
receives an \texttt{exit()} callback and the entering screen receives an
\texttt{enter()} callback. The exit callback can be used to store modified data
or release temporary resources, while the enter callback is typically used to
initialize the display contents and prepare the screen for user interaction.

A long press of the MENU button provides a convenient escape mechanism. It
returns the user to the first child of the root screen regardless of the current
position in the screen hierarchy.

\lstset{language=c++, style=codesnippetstyle}
\begin{lstlisting}
a conceptual picture of a screen hierarchy ?

MENU -> MENU -> MENU -> MENU -> MENU -> back to first menu
  |
  v
SUB-MENU
   |
   v
SUB-MENU
   |
   v
SUB-MENU
  |
  v
  back to first sub menu
\end{lstlisting}
\FloatBarrier

The default meaning of the MENU and SELECT buttons is a convention rather than
a restriction imposed by the framework. Individual screens may override the
corresponding event handlers and assign a different meaning whenever required.

All other UI elements are application defined. For example, a screen displaying
a value adjusted by a rotary encoder can override the encoder callback and react
to position changes directly. Likewise, additional buttons can be assigned
application-specific functions.

The connection between UI Elements and screens is established through the event
callback mechanism introduced in the previous chapter. UI Elements detect user
actions and forward the resulting events to the active screen. The screen then
decides how these events affect the application state.

The following section introduces the \texttt{UIScreen} object. It provides the
framework required to create application specific screens while connecting the
screen hierarchy to the underlying UI Elements. The actual implementation of
screens and complete user interfaces is discussed in the later chapter
\textbf{Programming with Screens}.

%--------------------------------------------------------------------------------
\unnumberedSection{The Screen Object}

The \texttt{UIScreen} object is the foundation for implementing application
specific screens. It combines two different responsibilities.

First, it manages the screen hierarchy. Each screen can have a parent screen and
a list of child screens. The base class provides the methods required to insert
screens into this hierarchy, enable or disable screens and determine the current
navigation position.

Second, it provides the connection between the UI Element library and the screen
application. UI Elements such as buttons and rotary encoders operate
independently of the screen hierarchy. During setup, the screen subsystem
registers callback functions with the required UI Elements. These static
callback functions receive the events generated by the UI Elements and forward
them to the currently active screen.

This separation is important because UI Elements represent the physical user
interface components, while screens represent the logical application behavior.
A button does not know which screen uses it, and a screen does not know whether
a button is connected directly to a controller pin or through another hardware
interface.

The \texttt{UIScreen} base class provides a number of virtual methods that are
intended to be overridden by application specific screens. These methods
implement the actual behavior of a screen, such as preparing the display when
entering a screen, reacting to button events or processing rotary encoder
movements.

The following listing shows the main structure of the \texttt{UIScreen} class.
The complete use of these methods is demonstrated in the chapter
\textbf{Programming with Screens}.
\lstset{language=c++, style=codesnippetstyle}
\begin{lstlisting}
struct UIScreen {

  public:

    UIScreen( );

    void            append( UIScreen *screen );

    virtual void    enterScreen( bool init );
    virtual void    exitScreen( );

    virtual void    menuButtonClick( UIButton *buttonId );
    virtual void    menuButtonLongPress( UIButton *buttonId );
    virtual void    selectButtonClick( UIButton *buttonId );
    virtual void    upButtonClick( UIButton *buttonId );
    virtual void    downButtonClick( UIButton *buttonId );

    virtual void    buttonClick( UIButton *buttonId );
    virtual void    buttonLongPressStart( UIButton *buttonId );
    virtual void    buttonLongPressStop( UIButton *buttonId );
    virtual void    encoderPosChange( UIEncoder *encoderObj );

    UIScreen        *getParentScreen( );
    UIScreen        *getChildScreen( );
    void            enableScreen( );
    void            disableScreen( );
    bool            isEnabled( );

    static void     menuButtonClickHandler( UIButton *buttonObj );
    static void     menuButtonLongPressHandler( UIButton *buttonObj );
    static void     selectButtonClickHandler( UIButton *buttonObj );
    static void     upButtonClickHandler( UIButton *buttonObj );
    static void     downButtonClickHandler( UIButton *buttonObj );

    static void     buttonClickHandler( UIButton *buttonObj );
    static void     buttonLongPressStartHandler( UIButton *buttonObj );
    static void     buttonLongPressStopHandler( UIButton *buttonObj );
    static void     encoderPosChangeHandler( UIEncoder *encoderObj );

    static bool     setup( );
    static UIScreen *getRootScreen( );
    static UIScreen *getCurrentScreen( );
    static UIScreen *setCurrentScreen( UIScreen *screen, bool init = true );
};

\end{lstlisting}
\FloatBarrier

The static callback methods at the end of the class definition form the bridge
between the UI Element library and the screen hierarchy. They are not
application callbacks. Their purpose is to receive hardware related UI events
and dispatch them to the active screen object.

The virtual methods, in contrast, define the application interface. A derived
screen class overrides only the methods relevant for its particular behavior.
For example, a screen displaying a value controlled by a rotary encoder would
typically override the \texttt{encoderPosChange()} method, while a configuration
screen may override button click handlers.

This design keeps the UI framework independent of individual applications while
allowing each application screen to implement exactly the required behavior.

%--------------------------------------------------------------------------------
\unnumberedSection{Summary}

The UI Screen library builds upon the UI Elements library by organizing
individual interface elements into complete operator interfaces. Rather than
processing button and encoder events throughout the application, user
interaction is concentrated in a hierarchy of screens that communicate through a
well-defined callback mechanism.

The \texttt{UIScreen} object provides the framework required to build these
interfaces. Static callback functions connect the screen subsystem to the UI
Elements, while virtual methods allow application specific screens to implement
their own behavior.

By separating navigation, event dispatching and application logic, the framework
encourages a consistent user interface while remaining sufficiently flexible for
specialized interaction patterns. Since screens operate exclusively on the
abstract UI elements introduced in the previous chapter, complete user
interfaces remain independent of the underlying controller hardware and display
technology.

Together, the UI Elements and UI Screen libraries provide the complete software
foundation for implementing portable and reusable operator interfaces across
the various LCS hardware modules. The practical implementation of screens and
complete user interfaces is presented when we implement cap handhelds, since
most of the UI elements screen requirements originates from there.
